%% LyX 2.4.4 created this file.  For more info, see https://www.lyx.org/.
%% Do not edit unless you really know what you are doing.
\documentclass[oneside,english]{amsart}
\usepackage[T1]{fontenc}
\usepackage[latin9]{inputenc}
\usepackage{units}
\usepackage{enumitem}
\usepackage{amsthm}

\makeatletter
%%%%%%%%%%%%%%%%%%%%%%%%%%%%%% Textclass specific LaTeX commands.
\numberwithin{equation}{section}
\numberwithin{figure}{section}
\newlength{\lyxlabelwidth}      % auxiliary length 

%%%%%%%%%%%%%%%%%%%%%%%%%%%%%% User specified LaTeX commands.
\usepackage{enumitem}
\usepackage{ifthen}
\usepackage{lastpage}
%sets math fonts to be more readable
%\everymath=\expandafter{\the\everymath\displaystyle}
%\DeclareMathSizes{10}{10}{10}{10}
\usepackage{multicol}
\usepackage{fancyhdr}
\pagestyle{fancy}
\usepackage{titlepic}
\usepackage{graphicx}
\usepackage{lineno}
\usepackage{hyperref}
%allows header and footer on titlepage by redefining the plain pagestyle issued by \maketitle to be fancy
%\makeatletter
%\let\ps@plain\ps@fancy
%\makeatother
\usepackage{advdate}    % Advancing/saving dates
\usepackage{datetime}   % Dates formatting
\usepackage{datenumber} % Counters for dates

\usepackage{hyperref}
\usepackage[usenames,dvipsnames]{xcolor} %adds additional color options
\hypersetup{colorlinks=true, urlcolor=Blue, linkcolor=Blue, citecolor=Blue}
\usepackage{url}
% Added by lyx2lyx
\setlength{\parskip}{\medskipamount}
\setlength{\parindent}{0pt}

\makeatother

\usepackage{babel}
\begin{document}
\renewcommand{\author}{Dr.\,James Ashe}

\newcommand{\classNum}{137}

\newcommand{\classSec}{all sections}

\newcommand{\nameType}{Math Lab Project}

\newcommand{\examType}{}

\renewcommand{\date}{}

%%First page's header%%%%%%%%%%%%%%%%%%%%%%
\fancypagestyle{firststyle} %%header settings for first pagr
{
	\fancyhead[L]{MTH \classNum{}(\classSec{}) \author{} \textbf{\examType{}} \date{}}
	\fancyhead[C]{}
	\fancyhead[R]{\textbf{\nameType}}
	\setlength{\headsep}{5pt}
}
%%%%%%%%%%%%%%%%%%%%%%%%%%%%%%%%%%%%%%%%%%

%%Next page's header%%%%%%%%%%%%%%%%%%%%%%%
\setlist{nolistsep}
\lhead{\classNum{}(\classSec{})} 
\chead{\textbf{\nameType}} 
\cfoot{\thepage{}~of~\pageref{LastPage}} 
\setlength{\headsep}{5pt} 
%%%%%%%%%%%%%%%%%%%%%%%%%%%%%%%%%%%%%%%%%%%

%%Various settings; see the preamble for other settings%%%%%
%%\setenumerate[0]{leftmargin=12pt,itemindent=0pt,labelwidth=0pt}
\setlist{itemsep=2pt}
\renewcommand{\labelenumi}{\textbf{\arabic{enumi}.}}
\renewcommand{\labelenumii}{\textbf{\alph{enumii}.}}
\renewcommand{\labelenumiii}{\textbf{\roman{enumiii}.}}
\thispagestyle{firststyle}
%%%%%%%%%%%%%%%%%%%%%%%%%%%%%%%%%%%%%%%%%%

\*

\section*{Completing the Square}
\begin{enumerate}
\item Using cardboard or poster board, construct a square.

\begin{enumerate}
\item The square should be large enough for a presentation.
\item Label each side $x$. 
\item Write $x^{2}$ in the center. This is the area of the square.\vfill
\end{enumerate}
\item Construct a rectangle with length $x$ (same as the square from step
1) and width $b$. 

\begin{enumerate}
\item $b$ can be any width you like, but it would be best if $b$ was between
a third to one-half of $x$.
\item Write $bx$ in the center of the rectangle. This is the area of the
rectangle.\vfill
\end{enumerate}
\item Explain how the square and rectangle represent $x^{2}+bx$.\vfill
\item Make an exact copy of the rectangle ``$bx$'' (from step 2).

\begin{enumerate}
\item Cut this rectangle in half.
\item Label the sides of both of these rectangles with $\nicefrac{b}{2}$
and $x$ appropriately. \vfill
\end{enumerate}
\item Arrange the ``$x^{2}$'' square and the two ``$\nicefrac{b}{2}$''
rectangles as follows:

\begin{enumerate}
\item Place one of the ``$\nicefrac{b}{2}$'' rectangles so that its side
of length $x$ is paired with a side of the ``$x^{2}$'' square
(this will form a rectangle).
\item Next place the other ``$\nicefrac{b}{2}$'' rectangle so that it
is orthogonal to the other ``$\nicefrac{b}{2}$'' rectangle and
also paired with a side of the ``$x^{2}$'' square (this will look
like a big square with a little square corner missing).\vfill
\end{enumerate}
\item Make a little square that will fit into the missing corner in the
above step. 

\begin{enumerate}
\item Place this little square into the above step to ``complete the square''.
The result should be a big square.
\item This little square has an area of $\left(\nicefrac{b}{2}\right)^{2}$;
explain.\vfill
\end{enumerate}
\item Determine the area of the big square from the above step.\vfill
\item Explain how the steps above relate to the equation: $x^{2}+bx+\left(\nicefrac{b}{2}\right)^{2}=\left(x+\nicefrac{b}{2}\right)^{2}$.\vfill
\item Relate how the concepts in the above steps can be used to write $x^{2}+6x+7$
in the form $\left(x+3\right)^{2}-2$.\vfill
\item Write an essay or prepare a presentation explaining the steps above.
Show all your work and include the shapes from steps 1-4.\vfill
\end{enumerate}

\end{document}
